%% ============================================================
%% VIKOR 段落模板
%% 变量:
%%   n_criteria, n_samples, criteria_names, object_names
%%   weights, weights_source
%%   S_values     — 群体效用值列表
%%   R_values     — 个体遗憾值列表
%%   Q_values     — 折中排序值列表
%%   v_param      — 决策机制系数 v (默认0.5)
%%   vikor_ranks  — VIKOR Q值排名
%%   result_bar_path — 结果柱状图路径
%% ============================================================

\subsection{基于 VIKOR 的折中排序模型}

\subsubsection{模型原理}

VIKOR（VlseKriterijumska Optimizacija I Kompromisno Resenje）方法由 Opricovic 于 1998 年提出，其核心理念是在群体效用最大化与个体遗憾最小化之间寻求折中解。

设有 $m = << n_samples >>$ 个评价对象、$n = << n_criteria >>$ 个指标，权重为 $\mathbf{w}$（由<< weights_source >>确定）。

\textbf{步骤一：确定正理想值 $f_j^*$ 和负理想值 $f_j^-$。}
\begin{equation}
    f_j^* = \max_i f_{ij}, \quad f_j^- = \min_i f_{ij} \quad \text{（正向指标）}
\end{equation}

\textbf{步骤二：计算群体效用 $S_i$ 和个体遗憾 $R_i$。}
\begin{equation}
    S_i = \sum_{j=1}^{n} w_j \frac{f_j^* - f_{ij}}{f_j^* - f_j^-}, \quad
    R_i = \max_j \left[ w_j \frac{f_j^* - f_{ij}}{f_j^* - f_j^-} \right]
\end{equation}

\textbf{步骤三：计算折中排序值 $Q_i$。}
\begin{equation}
    Q_i = v \frac{S_i - S^*}{S^- - S^*} + (1 - v) \frac{R_i - R^*}{R^- - R^*}
    \label{eq:vikor_Q}
\end{equation}
其中 $v = << v_param >>$ 为决策机制系数，$S^* = \min S_i$，$S^- = \max S_i$，$R^* = \min R_i$，$R^- = \max R_i$。$Q_i$ 越小表示综合表现越优。

\subsubsection{计算结果}

\begin{table}[H]
    \centering
    \caption{VIKOR 综合评价结果}
    \label{tab:vikor_result}
    \begin{tabular}{lcccc}
        \toprule
        \textbf{评价对象} & $S_i$ & $R_i$ & $Q_i$ & \textbf{排名} \\
        \midrule
<% for i in range(n_samples) %>
        << object_names[i] >> & << "%.6f" | format(S_values[i]) >> & << "%.6f" | format(R_values[i]) >> & << "%.6f" | format(Q_values[i]) >> & << vikor_ranks[i] >> \\
<% endfor %>
        \bottomrule
    \end{tabular}
\end{table>

<% if result_bar_path %>
\begin{figure}[H]
    \centering
    \includegraphics[width=0.8\textwidth]{<< result_bar_path >>}
    \caption{VIKOR $S$、$R$、$Q$ 值对比}
    \label{fig:vikor_result}
\end{figure}
<% endif %>

\subsubsection{折中解检验}

VIKOR 方法要求最终折中解同时满足"可接受优势"和"可接受稳定性"两个条件：
\begin{itemize}
    \item \textbf{可接受优势：} $Q(a'') - Q(a') \geq 1/(m-1)$，其中 $a'$ 和 $a''$ 分别为 $Q$ 值排名第一和第二的方案。
    \item \textbf{可接受稳定性：} $a'$ 在 $S$ 或 $R$ 排名中也为第一。
\end{itemize}
若两个条件均满足，则 $Q$ 值最小的方案为折中最优解。